\section{MuP 的实证验证}

\subsection{大规模 Ablation 研究}

\begin{figure}[H]
\centering
\includegraphics[width=0.95\textwidth]{slides-images/slide-040.jpg}
\caption{大规模 MuP 探索论文概述：系统验证 MuP 在各种架构变体、优化器和正则化选择下的鲁棒性\protect\footnotemark}
\end{figure}
\footnotetext{Slides 第41页。}

讲者介绍了一篇发表于 CoLM 的论文《A Large-Scale Exploration of $\mu$-Transfer》，该论文通过大量 ablation 实验系统检验了 MuP 的适用范围和局限性。

\subsection{基础验证：MuP 有效吗？}

\begin{figure}[H]
\centering
\includegraphics[width=0.95\textwidth]{slides-images/slide-041.jpg}
\caption{MuP 有效性验证：不同宽度（128--2048）下，最优学习率（$2^{-6}$）保持不变\protect\footnotemark}
\end{figure}
\footnotetext{Slides 第42页。}

实验设置：固定深度，将宽度从 128 扩展到 2048。在最小宽度上扫描学习率，找到最优值。如果 MuP 有效，这个最优学习率应该直接迁移到更大宽度。

结果：\textbf{学习率确实可靠迁移}。不同宽度的最优学习率都是 $2^{-6}$。

\subsection{MuP 对什么鲁棒？}

\begin{figure}[H]
\centering
\includegraphics[width=0.95\textwidth]{slides-images/slide-042.jpg}
\caption{非线性变体：SwiGLU、Squared ReLU 等都不影响最优学习率的迁移\protect\footnotemark}
\end{figure}
\footnotetext{Slides 第43页。}

论文测试了多种架构变体：

\begin{table}[H]
\centering
\begin{tabular}{lcc}
\toprule
\textbf{变体} & \textbf{学习率迁移} & \textbf{备注} \\
\midrule
非线性激活（SwiGLU、Squared ReLU） & 有效 & 最优 LR 不变 \\
Batch size 变化（$\times 4$ 上下） & 有效 & 最优 LR 不变 \\
初始化变体（Query 初始化为零等） & 有效 & 最优 LR 不变 \\
Unembedding 层 SP vs MuP & 有效 & 最优 LR 不变 \\
\bottomrule
\end{tabular}
\caption{MuP 鲁棒性测试——有效的情况}
\end{table}

\begin{figure}[H]
\centering
\includegraphics[width=0.95\textwidth]{slides-images/slide-043.jpg}
\caption{Batch size 变体、初始化变体的学习率迁移测试：均保持稳定\protect\footnotemark}
\end{figure}
\footnotetext{Slides 第44页。}

\subsection{MuP 在何时失效？}

\begin{figure}[H]
\centering
\includegraphics[width=0.95\textwidth]{slides-images/slide-045.jpg}
\caption{MuP 失效案例：可学习的 RMSNorm gain（左上）、Lion 优化器（右上）、强 weight decay（下方）\protect\footnotemark}
\end{figure}
\footnotetext{Slides 第46页。}

\begin{warningbox}{MuP 失效的三种情况}
\begin{enumerate}
    \item \textbf{可学习的 bias/gain}：如果在 RMSNorm 中添加可学习的 gain 参数，MuP 会失效。需要移除这些参数。
    \item \textbf{非标准优化器}：如 Lion（取梯度的 sign 作为更新方向）。MuP 是为特定优化器（SGD/Adam/AdamW）设计的，使用完全不同的优化器自然无法保证学习率迁移。
    \item \textbf{强 Weight Decay}：较强的 weight decay 会破坏 MuP 的学习率迁移。这是一个显著的实践限制，因为 weight decay 是标准训练中常用的正则化手段。
\end{enumerate}
\end{warningbox}

\subsection{标准参数化 vs MuP 的对比}

\begin{figure}[H]
\centering
\includegraphics[width=0.95\textwidth]{slides-images/slide-047.jpg}
\caption{标准参数化（SP）的学习率迁移：相同学习率在宽度 2048 时导致模型爆炸/退化\protect\footnotemark}
\end{figure}
\footnotetext{Slides 第48页。}

使用标准参数化时，在小模型上找到的最优学习率\textbf{不能}直接用于大模型——大模型会因为更新过大而训练失败。学习率必须随规模\textbf{可预测地下降}。

\subsection{大规模验证}

\begin{figure}[H]
\centering
\includegraphics[width=0.95\textwidth]{slides-images/slide-048.jpg}
\caption{10B 参数模型验证：在中小规模上确定的 base 学习率 $2^{-6}$ 在 10B 模型上仍然是最优的\protect\footnotemark}
\end{figure}
\footnotetext{Slides 第49页。}

论文还进行了一次大规模``hero run''验证：将在中小规模上找到的最优学习率直接用于 10B 参数模型，结果证实该学习率仍然是最优的。这为 MuP 的实用性提供了重要的经验支持。

\begin{knowledgebox}{Meta 的 Llama 4 与 MuP}
讲者提到，Meta 在 Llama 4 中使用了一种名为 ``MetaP'' 的技术，这是 MuP 的一个变体。虽然 Llama 4 的论文尚未发表，但这表明 MuP 类方法正在被前沿实验室采用。不过，MuP 目前还不是行业共识——并非所有团队都使用它。
\end{knowledgebox}

\subsection{本章小结}

经验验证表明，MuP 在大多数标准设置下（Adam/AdamW 优化器、常见非线性激活、合理的 batch size 范围）可以可靠地实现学习率的跨规模迁移。主要失效场景包括可学习的 norm gain、非标准优化器和强 weight decay。

%% ========== Part 8: 总结 ==========

